% !TEX root = ../metatime_monograph.tex
\chapter{Publication Protocol, Data Provenance, and Statistical Interpretation}
\label{app:publication_protocol}

% CITATION_CONTEXT_V10_9_BEGIN
\paragraph{Established context.}
Model comparison, likelihood inference, preregistration, and reproducible
computation are grounded in the statistical and methodological literature
\cite{Akaike1974,Schwarz1978,Wilks1938,Cowan2011,Nosek2018,Sandve2013}.
The protocol below applies those norms to interpretation of this monograph.
% CITATION_CONTEXT_V10_9_END

\claimstatus{Methodological protocol; it governs interpretation but is not a
physical derivation.}

\section{Purpose}
This appendix fixes the publication-level interpretation of the numerical
material in the monograph.  It prevents a close numerical match from being
promoted to independent evidence when the corresponding formula was selected
retrospectively, and it prevents quantities with incompatible statistical
meanings from being merged into a single accuracy score.  The protocol follows
standard concerns in model selection, likelihood inference, preregistration,
and reproducible computation
\cite{Akaike1974,Schwarz1978,Wilks1938,Cowan2011,Nosek2018,Sandve2013}.

\section{Data snapshot and external inputs}
The numerical comparison tables are a frozen audit snapshot dated 29 July 2026.
Unless an entry explicitly states otherwise, a quoted reference value is the
value used during development of the v10 series and is not asserted to be the
latest global fit.  The principal external reference compilations are PDG,
CODATA, and Planck \cite{PDG2024,CODATA2022,Planck2018}.

External inputs include the Planck scale, the proton mass where used as a
hadronic scale, $N_c=3$, the assigned lepton hypercharge, and the hadronic
coefficient used to map $\theta_{\rm QCD}$ to the neutron EDM.  The electron is
the sole dimensional anchor in the Collatz charged-fermion mass audit.  Results
depending on an anchor are not independent predictions of that anchor.

\section{Renormalization and scheme conventions}
The current monograph does not derive renormalization-group running.  Therefore:
\begin{enumerate}
  \item masses quoted for charged leptons, hadrons, and gauge bosons are compared
  with the reference mass conventions stated or implied by the cited source;
  \item $\sin^2\theta_W$, gauge couplings, quark masses, and Yukawa couplings are
  scheme- and scale-dependent and must not be compared at precision level unless
  the same scheme and scale are used;
  \item future ratios $y_c/y_\mu$ and $y_c/y_t$ must be taken from a common direct
  experimental likelihood or translated to a common scheme before testing the
  frozen candidates;
  \item no sub-percent claim is made for a cross-scheme comparison.
\end{enumerate}

\section{Residuals and aggregate summaries}
For a positive observable $x$, the preferred scale-free residual is
\[
  r=\ln\!\left(\frac{x_{\rm model}}{x_{\rm ref}}\right),
  \qquad
  f=\exp(|r|),
\]
where $f$ is the multiplicative error factor.  A normalized pull may be quoted
only when a defensible model uncertainty and reference uncertainty are both
available:
\[
  z=\frac{x_{\rm model}-x_{\rm ref}}
  {\sqrt{\sigma_{\rm model}^2+\sigma_{\rm ref}^2}}.
\]
Circular variables such as phases require circular distance.  Upper limits such
as the neutron EDM are pass/fail constraints and are excluded from mean-error
calculations.  External anchors and values used to choose a formula are also
excluded from prospective performance scores.

The earlier statement of a mean error across ``26 SM parameters'' is withdrawn.
The summary actually contains 36 heterogeneous observables and derived
quantities, including hadron masses and cosmological quantities that are not
independent parameters of the Standard Model Lagrangian.  Sector-level
descriptive residuals may be reported, but no single global percentage is used
as evidence.

\section{Model complexity}
``Zero continuous fitted parameters'' does not mean zero model complexity.  The
construction contains approximately fifty discrete structural choices,
including integer assignments, prime labels, formula families, anchor choices,
and sector mappings.  Interpretation must therefore account for discrete
selection freedom and post-selection.  Information criteria such as AIC or BIC
are not applied here because a complete likelihood and effective parameter count
have not yet been defined; their absence is a declared limitation, not evidence
in favor of the model.

\section{Retrospective and prospective evidence}
All numerical sector formulas developed with access to their target values are
classified as retrospective assignments.  They may motivate structure but do
not constitute independent confirmation.  The v10.7 candidate family is the
principal prospective component because the three formulas, two orthogonal
observables, post-29-July-2026 data gate, and no-refit rule are frozen in advance.
Multiplicity-aware likelihood interpretation is mandatory.

\section{Current physical failures}
\begin{enumerate}
  \item The fixed gauge-boson relations overshoot the quoted $W$ and $Z$ masses
  by approximately $4.5\%$ and $5.0\%$.
  \item With $\kappa=\ln2/(24\pi)$ and exponent $14$, the strong-CP assignment
  gives $\theta_{\rm QCD}=2.2208\times10^{-10}$ and, with the fixed hadronic
  coefficient used here, $d_n=5.3299\times10^{-26}\,e\,\mathrm{cm}$.  This is a
  technical calculation PASS but a physical constraint FAIL against the
  $1.8\times10^{-26}\,e\,\mathrm{cm}$ bound used in the manuscript.
  \item The isolated Collatz $3/4$ mass trace has a geometric-mean multiplicative
  error of $9.967$ and is not a closed spectrum derivation.
\end{enumerate}

These failures are retained without hidden suppression factors, exponent changes,
or candidate substitution.

\section{Publication gate}
A release is a publication candidate only if the build has no unresolved
citations or references, no multiply defined labels, no overfull boxes, nonempty
PDF metadata, embedded non-Type-3 fonts, a complete citation-coverage ledger, and
a publication-readiness ledger recording the physical PASS/FAIL separation.
